\section{Algorithm}

The GIA ANN prototype is a concept--feature neural database that learns from text sequences and predicts next tokens by activating connected features. It treats nouns as concept columns, and every other token as a feature neuron within those columns. The algorithm is designed for online, unbatched learning, and stores most of its state in sparse tensors for scalability.

\begin{center}
\includesvg[width=\linewidth]{demo/GIAANNproto1cAllColumnsTrainSequenceIndex005}
\par\smallskip
\refstepcounter{figure}\label{fig:giaannproto1c_allcolumns_train_005}
\textbf{Figure~\thefigure:} Example all-columns network visualisation during training.
\end{center}

\subsection{Core Data Structures}
\begin{itemize}
	\item \textbf{Concept columns:} a dynamic dictionary mapping concept lemmas to column indices ($c$). Each column has a prime concept neuron at a fixed feature index.
	\item \textbf{Feature dictionary:} a dynamic dictionary mapping feature words to feature indices ($f$). Features are words (not lemmas) and are shared across columns.
	\item \textbf{Observed columns:} per-sequence objects that hold the sparse feature neuron tensors and feature connection tensors for each concept column seen in the sequence.
	\item \textbf{SequenceObservedColumns:} a per-sequence view that stacks the observed columns into dense or sparse tensors of shape (properties$\times$branches$\times$segments$\times$columns$\times$features), plus a matching feature--connection tensor.
	\item \textbf{Properties:} strength, permanence, activation, time, POS, and optional minimum word distance. Which properties exist is controlled by configuration.
	\item \textbf{SANI segments:} each dendritic branch is segmented to encode sequential ordering. Segment indices are assigned by column distance, feature distance, or both.
\end{itemize}

\subsection{Preprocessing and Column Segmentation}
\begin{enumerate}
	\item \textbf{Sequence construction:} text is parsed with spaCy; sequences are either single sentences or multi-sentence bundles. Empty or too-short sequences are skipped.
	\item \textbf{Token normalisation:} tokens are lowercased, lemmatised, and optionally preprocessed to combine consecutive nouns or hyphenated noun phrases.
	\item \textbf{Concept detection:} nouns are detected via POS rules (or the full POS dictionary); these lemmas define concept columns.
	\item \textbf{Delimiter detection:} reference set delimiters are detected by POS, tag, or word lists. Deterministic delimiters take priority over probabilistic ones.
	\item \textbf{Column assignment:} tokens are assigned to concept columns by choosing the rightmost delimiter between adjacent concept tokens. If no delimiter exists, the sequence is skipped (training) or inference is aborted (prediction).
\end{enumerate}

\subsection{Training Algorithm}
\begin{enumerate}
	\item \textbf{First pass (concept growth):} scan the sequence and add new concept lemmas to the concept column dictionary. If $c$ grows, resize global feature tensors.
	\item \textbf{Feature discovery:} scan the sequence for non-concept tokens and add new feature words to the feature dictionary. If $f$ grows, expand all feature tensors.
	\item \textbf{Observed column loading:} for each concept in the sequence, load or create its observed column object and refresh its arrays to the current $c$ and $f$.
	\item \textbf{Column segmentation:} compute per-concept token spans and build a token-to-column assignment list.
	\item \textbf{Feature activation tensor:} for each concept instance, mark active feature neurons inside the column span. Assign dendritic branches (optionally randomized) and SANI segments based on column or feature distance.
	\item \textbf{Connection activation tensor:} create directed connections between active features, enforcing forward word order and column order. Optionally enforce direct connections (adjacent tokens only).
	\item \textbf{Update rule:}
	\begin{itemize}
		\item Strength is incremented for active features and connections (with optional normalisation and POS-dependent scaling).
		\item Permanence is incremented for active features/connections and can be decremented for inactive ones.
		\item Activation traces and time properties are reset or updated as configured.
		\item POS and minimum word distance properties are recorded for active items.
	\end{itemize}
	\item \textbf{Commit:} merge sequence tensors back into observed columns, save updated columns and dictionaries to disk, and optionally draw the network.
\end{enumerate}

\subsection{Inference Algorithm (Seed + Prediction)}
\begin{enumerate}
	\item \textbf{Seed phase:} use the same prediction pipeline as normal inference, but force the selected feature to be the target token. This seeds the activation state.
	\item \textbf{Activation state:} use global sparse activation tensors (and optional connection tensors) as the working state for predictions.
	\item \textbf{Sequential gating (SANI):} apply sequential activation rules so a segment can only fire if the previous segment is active.
	\item \textbf{Delimiter constraints:} if the current active node is a reference-set delimiter, constrain the next prediction to delimiter-compatible nodes or external columns.
	\item \textbf{Time-based biasing:} optionally adjust segment activations based on last activation time relative to the expected segment index.
	\item \textbf{Candidate selection:}
	\begin{itemize}
		\item If beam search is enabled, expand beams to a fixed depth and width, scoring by activation and connection strength.
		\item Otherwise, select the single best node from the current activation map.
	\end{itemize}
	\item \textbf{State update:} activate the chosen node, update activation traces and time properties, and propagate its outgoing connections to update the next step.
	\item \textbf{Evaluation:} optionally compare predictions to the target token and log mismatches.
\end{enumerate}

\subsection{Beam Search (Optional)}
\begin{itemize}
	\item Each beam holds a cloned activation state (feature activations, optional connection activations, and time properties).
	\item Candidate nodes are filtered by constraint mode, connected-column constraints, and activation thresholds.
	\item Beam scoring can use activation gain, connection strength, or both. The final prediction is the first action of the highest-scoring beam.
\end{itemize}

\subsection{Persistence and Visualisation}
\begin{itemize}
	\item Each observed column is stored on disk with its sparse tensors and feature mappings.
	\item Global feature neuron tensors and concept/feature dictionaries are saved at configurable intervals.
	\item The visualiser renders concept columns, feature nodes, and connections with colour schemes for segments, branches, POS types, or delimiters.
\end{itemize}
